\section*{Capítulo \ref{ch:g10:habitable-planet} --- Um planeta habitável}

\begin{solution}{exo:g10:habitable-planet:1}
A temperatura em que um emissor perfeito irradiaria para o espaço
exatamente a potência que o planeta absorve da sua estrela. A superfície da
Terra é mais quente (\qty{288}{K} contra cerca de \qty{255}{K}) porque os
gases de efeito estufa absorvem parte do infravermelho que ela emite e
devolvem parte dele para baixo.
\end{solution}

\begin{solution}{exo:g10:habitable-planet:2}
Vênus: \qty{737}{K} medidos contra uma \omterm{def:g10:habitable-planet:equilibrium-temperature}{temperatura de equilíbrio} de cerca de
\qty{227}{K}, um aquecimento de cerca de \qty{510}{K}.
\end{solution}

\begin{solution}{exo:g10:habitable-planet:3}
$S = 1361/1.52^2 \approx \qty{589}{W/m^2}$ (a ficha técnica dá
\qty{586}{W/m^2}, com a distância média exata).
\end{solution}

\begin{solution}{exo:g10:habitable-planet:4}
\qty{1}{bar}, \qty{330}{K}: líquido (abaixo do ponto de ebulição, cerca de
\qty{373}{K}). \qty{0.001}{bar}, \qty{250}{K}: gelo (a pressão fica acima
da curva gelo--vapor, que está em cerca de \qty{0.0008}{bar} a
\qty{250}{K}). \qty{100}{bar}, \qty{700}{K}: vapor --- a temperatura está
acima da temperatura crítica, logo nenhum líquido pode existir.
\end{solution}

\begin{solution}{exo:g10:habitable-planet:5}
Um planeta que orbita uma estrela que não é o Sol. Trânsitos: o planeta
passa na frente da sua estrela e a escurece levemente a cada órbita.
Oscilação: a atração do planeta move a estrela para a frente e para trás ao
longo da nossa linha de visada, o que desloca periodicamente a cor da sua
luz.
\end{solution}

\begin{solution}{exo:g10:habitable-planet:6}
$T = (1361 \times 0.706/(4 \times 5.67\times10^{-8}))^{1/4} \approx
\qty{255}{K} \approx \qty{-18}{\degreeCelsius}$. Com $A = 0$:
$\approx \qty{278}{K} \approx \qty{5}{\degreeCelsius}$.
\end{solution}

\begin{solution}{exo:g10:habitable-planet:7}
$\sqrt 4 = 2$: de $2 \times 0.99 = \qty{1.98}{au}$ a $2 \times 1.70 =
\qty{3.40}{au}$. A \qty{1}{au}, o planeta receberia quatro vezes a
irradiância da Terra: muito aquém da borda interna, quente demais.
\end{solution}

\begin{solution}{exo:g10:habitable-planet:8}
$S = 1361/0.723^2 \approx \qty{2604}{W/m^2}$;
$T_{\mathrm{eq}} = (2604 \times 0.23/(4\sigma))^{1/4} \approx \qty{227}{K}$.
A superfície é cerca de \qty{510}{K} mais quente: a sua espessa atmosfera de
gás carbônico (\qty{92}{bar}) produz um efeito estufa enorme.
\end{solution}

\begin{solution}{exo:g10:habitable-planet:9}
A \qty{6.36}{mbar}, a água só é líquida entre cerca de
\qty{0}{\degreeCelsius} e \qty{0.5}{\degreeCelsius}. A
\qty{10}{\degreeCelsius}, a pressão fica abaixo da curva de ebulição: a
água se evapora fervendo; à noite, ela congela.
\end{solution}

\begin{solution}{exo:g10:habitable-planet:10}
$11.2/2.4 \approx 4.7$. Uma molécula de gás precisa se mover quase cinco
vezes mais depressa para escapar da Terra do que da Lua; ao longo de bilhões
de anos, a gravidade fraca da Lua deixou os seus gases escaparem, enquanto a
da Terra reteve a maior parte deles.
\end{solution}

\begin{solution}{exo:g10:habitable-planet:11}
$T = (1361 \times 0.50/(4\sigma))^{1/4} \approx \qty{234}{K}$, cerca de
\qty{-39}{\degreeCelsius}. O gelo refletiria tanta luz do Sol que o planeta
ficaria frio o bastante para conservá-lo: um ciclo que se reforça.
\end{solution}

\begin{solution}{exo:g10:habitable-planet:12}
$T_{\mathrm{eq}} = (S_0(1-A)/(4\sigma d^2))^{1/4} = (S_0/(4\sigma))^{1/4}
(1-A)^{1/4} (d^2)^{-1/4}$; $(1361/(4\sigma))^{1/4} \approx \qty{278.3}{K}$ e
$(d^2)^{1/4} = \sqrt d$. Quatro vezes mais longe: $T$ dividida por $\sqrt 4 = 2$.
\end{solution}

\begin{solution}{exo:g10:habitable-planet:13}
$\sqrt d = 278.3 \times 0.706^{1/4}/273 \approx 0.934$, logo $d \approx
\qty{0.873}{au}$, aquém da borda interna (\qty{0.99}{au}). A \omterm{def:g10:habitable-planet:equilibrium-temperature}{temperatura de
equilíbrio} ignora o efeito estufa: um planeta real com atmosfera é mais
quente que $T_{\mathrm{eq}}$, logo ele superaquece e perde a sua água já
mais longe, a cerca de \qty{0.99}{au}. As bordas vêm de modelos da
atmosfera inteira, não apenas de $T_{\mathrm{eq}}$.
\end{solution}

\begin{solution}{exo:g10:habitable-planet:14}
Entre o ponto de fusão, cerca de \qty{0}{\degreeCelsius}, e o ponto de
ebulição, cerca de \qty{7}{\degreeCelsius}. A \qty{15}{\degreeCelsius}, o
lago ferveria.
\end{solution}

\begin{solution}{exo:g10:habitable-planet:15}
$\sqrt{0.5} \approx 0.707$: de $0.707 \times 0.99 \approx \qty{0.70}{au}$
a $0.707 \times 1.70 \approx \qty{1.20}{au}$. O planeta está exatamente na
borda externa. As bordas são resultados de modelos e o ajuste é
aproximado, logo é impossível dizer se ele está dentro ou fora.
\end{solution}

\begin{solution}{pb:g10:habitable-planet:1}
\textbf{1.} $0.00151/0.0485^2 \approx 0.642$.

\textbf{2.} $0.642 \times 1361 \approx \qty{874}{W/m^2}$.

\textbf{3.} Entre as duas: cerca de três vezes menos que em Vênus, uma vez
e meia mais que em Marte.

\textbf{4.} $365.25/11.18 \approx 32.7$.

\textbf{5.} $(874/(4 \times 5.67\times10^{-8}))^{1/4} \approx \qty{249}{K}$.

\textbf{6.} $(874 \times 0.23/(4\sigma))^{1/4} \approx \qty{173}{K}$.

\textbf{7.} Nenhuma: \qty{249}{K} ainda fica abaixo de \qty{273}{K}. Uma
cobertura clara de nuvens ou de gelo baixa a \omterm{def:g10:habitable-planet:equilibrium-temperature}{temperatura de equilíbrio} em
mais de \qty{75}{K}: o que cobre um planeta importa tanto quanto a sua
distância.

\textbf{8.} $\sqrt{0.00151} \approx 0.0389$.

\textbf{9.} $0.99 \times 0.0389 \approx \qty{0.0385}{au}$.

\textbf{10.} $1.70 \times 0.0389 \approx \qty{0.066}{au}$.

\textbf{11.} Sim: $0.0385 < 0.0485 < 0.066$.

\textbf{12.} $d_\odot = 0.0485/0.0389 \approx \qty{1.25}{au}$: na metade
interna da \omterm{def:g10:habitable-planet:habitable-zone}{zona habitável} do Sol, entre a Terra e Marte.

\textbf{13.} $1.055^{1/3} \approx 1.018$.

\textbf{14.} $11.2 \times \sqrt{1.055/1.018} \approx \qty{11.4}{km/s}$,
cerca da velocidade da Terra: a sua gravidade poderia reter uma atmosfera
(se ele tem uma depende da sua história).

\textbf{15.} A oscilação da estrela revela apenas a parte do seu movimento
ao longo da nossa linha de visada; como a inclinação da órbita é
desconhecida, a massa encontrada só é a massa verdadeira se a órbita for
vista de perfil, e é um mínimo nos outros casos.

\textbf{16.} A pressão do ponto triplo, \qty{611.657}{Pa}, cerca de
\qty{0.006}{bar}.

\textbf{17.} $(874 \times 0.706/(4\sigma))^{1/4} \approx \qty{228}{K}$.

\textbf{18.} $288 - 228 = \qty{60}{K}$: Proxima~b precisaria de um
aquecimento por efeito estufa de \textbf{cerca de \qty{60}{K}}, quase o
dobro dos \qty{33}{K} da Terra.
\end{solution}
